% ============================================================
% Conditional Stationarity and Positive Curvature of the
% Spectral Trace at the Critical Line
%
% Author:   Ulrich Tehrani
% Date:     May 2026
% Version:  v0.14
% Zenodo:   [DOI to be assigned upon publication]
% License:  CC BY 4.0
% ============================================================

\documentclass[11pt]{article}
\usepackage{cmap}
\usepackage[T1]{fontenc}
\usepackage[utf8]{inputenc}
\usepackage[a4paper,margin=1in]{geometry}
\usepackage{amsmath,amssymb,amsthm,mathtools}
\usepackage{xcolor}
\usepackage{hyperref}
\usepackage{enumitem}
\usepackage{setspace}
\usepackage{booktabs}
\usepackage{array}
\usepackage{graphicx}

\hypersetup{colorlinks=true,
  linkcolor=blue!50!black,
  urlcolor=blue!50!black,
  citecolor=blue!50!black}
\setlength{\parskip}{0.5em}
\setlength{\parindent}{0pt}
\onehalfspacing
\emergencystretch=1em

% --------------------------------------------------------
% Theorem environments (numbered within section)
% --------------------------------------------------------
\newtheorem{theorem}{Theorem}[section]
\newtheorem{proposition}[theorem]{Proposition}
\newtheorem{lemma}[theorem]{Lemma}
\newtheorem{corollary}[theorem]{Corollary}
\newtheorem{definition}[theorem]{Definition}
\newtheorem{assumption}[theorem]{Assumption}
\newtheorem{observation}[theorem]{Observation}
\newtheorem{conjecture}[theorem]{Conjecture}
\newtheorem*{remark}{Remark}
\newtheorem{openproblem}[theorem]{Open Problem}

% Heuristic environment (clearly marked, never used in proofs)
\newenvironment{heuristic}[1]{%
  \par\smallskip\noindent\textit{[Heuristic: #1]}\quad}{%
  \par\smallskip}

% --------------------------------------------------------
% Convention box --- clean ruled box, no mdframed
% --------------------------------------------------------
\newenvironment{conventionbox}{%
  \medskip
  \noindent\rule{\linewidth}{0.4pt}\smallskip\par\noindent
}{%
  \par\smallskip\noindent\rule{\linewidth}{0.4pt}
  \medskip
}

% --------------------------------------------------------
% Series-wide macros
% --------------------------------------------------------
\newcommand{\Hstr}{H_{\mathrm{str}}}
\newcommand{\Hnull}{H_{\mathrm{null}}}
\newcommand{\Estr}{E_{\mathrm{str}}}
\newcommand{\Espec}{E_{\mathrm{spec}}}
\newcommand{\etaorig}{\eta_{\mathrm{orig}}}
\newcommand{\Tren}{T_{\mathrm{ren}}}
\newcommand{\Tt}{\widetilde{T}}
\newcommand{\Gun}{G^{\mathrm{un}}}
\newcommand{\Hlocal}{H_{\mathrm{local}}}
\newcommand{\DSEL}{D_{\mathrm{SEL}}}
\newcommand{\Ceta}{C_{\eta}}
\newcommand{\CT}{C_{T}}

% --------------------------------------------------------
% Paper-7-specific macros
% --------------------------------------------------------
\newcommand{\Wgamma}{\mathcal{W}_{\varepsilon,N}}   % Σ_k e^{-ε²γ_k²} γ_k (γ_k-weighted)
\newcommand{\Astar}{\ensuremath{(A^{**})}}  % Cancellation Hypothesis label

% --------------------------------------------------------
% Title
% --------------------------------------------------------
\title{\vspace{-1.5em}\bfseries Conditional Stationarity and Positive
Curvature of the Spectral Trace at the Critical Line}
\author{Ulrich Tehrani}
\date{Research Note \textperiodcentered\ May~2026
      \textperiodcentered\ v0.14}

% ============================================================
\begin{document}
\maketitle
% ============================================================

\begin{abstract}
Building on the curvature identity $O''(\tfrac12)=-2B$ established
in~\cite{Paper6} and the trace formula $\mathrm{Tr}(\Tt(\sigma))=\DSEL-O(\sigma)$
proved in~\cite{Paper5}, we address two of the five open problems
of Paper~6 concerning the first derivative $O'(\tfrac12)$ and the
pointwise asymptotic behaviour of the Guinand--Weil constant term.
Under an explicit subleading archimedean hypothesis
(Assumption~\ref{ass:subleading}, numerically supported) and a
proxy-transfer hypothesis (Assumption~\ref{ass:proxy},
formally weaker in content than RH: RH implies it, no converse claimed),
we prove that the ratio $r_p^\infty := \lim_{\varepsilon\to 0^+}
|\mathrm{Const}_p^\infty(\varepsilon)|/|\mathrm{Main}_p(\varepsilon)|$ equals
$\tfrac12$ for every prime $p$, and derive an exact
three-term bookkeeping decomposition of $O'(\tfrac12)$
with two explicit leading terms and a defined residual.
We further prove under the same assumptions that for every prime cutoff
$\kappa\ge 2$ there exists $\varepsilon_{\mathrm{int}}(\kappa)>0$ such that the
integrated bias $B_{\mathrm{int}}^\infty(\kappa,\varepsilon)<0$ for all
$\varepsilon\in(0,\varepsilon_{\mathrm{int}}(\kappa))$.
Under an additional logarithmic-saving hypothesis on prime exponential sums
(Assumption~\Astar, Cancellation Hypothesis) we establish a conditional
constant-factor bound on~$|O'(\tfrac12)|$ relative to the trivial bound
(Proposition~\ref{prop:deriv}).
The full conditional local-selection statement (strict local minimum of
$O$ at $\sigma=\tfrac12$) additionally requires the Weighted Bias Bridge
hypothesis $B(\kappa,\varepsilon,N)<0$ and exact stationarity
$O'(\tfrac12)=0$ (Corollary~\ref{cor:selection}).
No hypothetical input is used in the unconditional results: no Riemann
Hypothesis, no Hilbert--P\'olya postulate, no GUE conjecture.

\textbf{What is established unconditionally.}
Definition~\ref{def:threeterm} (bookkeeping decomposition);
Proposition~\ref{prop:antisym} (antisymmetry).
Observation~\ref{obs:threeterm} provides numerical values at
reference parameters and is reported under \emph{What is numerical}.

\textbf{What is proved under Assumptions~\ref{ass:subleading} and~\ref{ass:proxy}.}
For every prime $p$, $r_p^\infty = \tfrac12$ (Theorem~\ref{thm:rpinf});
in particular $Z_{p,\infty}^+(\varepsilon)<0$ for all sufficiently small
$\varepsilon>0$ (Corollary~\ref{cor:sign});
and for every $\kappa\ge 2$ there exists $\varepsilon_{\mathrm{int}}(\kappa)>0$
such that $B_{\mathrm{int}}^\infty(\kappa,\varepsilon)<0$ for
$\varepsilon\in(0,\varepsilon_{\mathrm{int}}(\kappa))$
(Proposition~\ref{prop:intbias}).

\textbf{What is numerical.}
At $(\kappa,\varepsilon,N)=(53,0.05,100)$:
$O'(\tfrac12)=+2.4751$; $O''(\tfrac12)=+38\,685.1$; $B=-19\,342.5$.
The ratio
$|O'(\tfrac12)|/(\mathcal{W}_{\varepsilon,N}\cdot P(\kappa))\approx 0.00194$
quantifies the cancellation relative to the trivial bound.
Separately, the internal balance among the three terms gives
$|T_{\mathrm{pol}}+T_{\mathrm{main}}|/|O'(\tfrac12)|\approx 15.7$
(see Observation~\ref{obs:threeterm} for the explicit values).
$\varepsilon_{\mathrm{int}}(53)>0.200$; $O''(\tfrac12)=+38\,685.1>0$ [numerical].

\textbf{Additional numerical diagnostics.}
Sections~\ref{sec:cancel}, \ref{sec:struct}, and~\ref{sec:open}
report further finite-grid numerical diagnostics, including the value
$O'(\tfrac12;101,0.05,100)=-137.6$ used in the structural-necessity
discussion, observed sign-region data for
$\varepsilon_{\mathrm{int}}(\kappa)$, and the tested near-minimum
behaviour at $\kappa=53$ and $\kappa=101$ (fixed $N=100$).
These are finite-grid numerical observations and are not used as
proved inputs.
The strong $\kappa$-dependence of $O'(\tfrac12;\kappa,\varepsilon,N)$
between $\kappa=53$ and $\kappa=101$ is precisely the empirical signature
that motivates introducing the cancellation hypothesis~\Astar:
exact stationarity at the reference parameters is not claimed to be
universal, and the variation across $\kappa$ is the structural reason
why an analytical input on prime-exponential-sum cancellation is
required to control $|O'(\tfrac12)|$ uniformly.

\textbf{What is conditional.}
Under Assumption~\Astar\ (Cancellation Hypothesis, \S\,\ref{sec:cancel}):
a constant-factor bound on $|O'(\tfrac12;\kappa,\varepsilon,N)|$ relative
to the trivial bound (Proposition~\ref{prop:deriv}).
Under the additional hypotheses $B(\kappa,\varepsilon,N)<0$ and
$O'(\tfrac12)=0$ (exact stationarity), this gives a conditional
local-selection criterion at $\sigma=\tfrac12$
(Corollary~\ref{cor:selection}).

\textbf{What remains open.}
The Weighted Bias Bridge --- whether $B_{\mathrm{int}}^\infty<0$ implies
$O''(\tfrac12)>0$ via $O''(\tfrac12)=-2B$ (Open Problem~\ref{op:wbridge});
the Cancellation Hypothesis~\Astar\ (Open Problem~\ref{op:cancel});
the Guinand--Weil bridge with off-critical control
(Open Problem~\ref{op:gwbridge}); the scaling of $\varepsilon_{\mathrm{int}}(\kappa)$
(Open Problem~\ref{op:epsc}); the near-minimum question
(Open Problem~\ref{op:global}); and the Weil-transfer connection
(Open Problem~\ref{op:weil}).

\textbf{Structure.}
The unconditional part (Definition~\ref{def:threeterm}, Proposition~\ref{prop:antisym})
is proof-complete and non-circular.
Results conditional on Assumptions~\ref{ass:subleading}
and~\ref{ass:proxy} are clearly labelled.
No use of the Riemann Hypothesis, the GUE conjecture, or the
Hilbert--P\'olya postulate is made in \S\S\,\ref{sec:bg}--\ref{sec:curv};
Assumption~\Astar\ in \S\,\ref{sec:sel} is explicitly declared conditional.
\end{abstract}

\vspace{2em}

% --------------------------------------------------------
% Notation and Conventions
% --------------------------------------------------------
\section*{Notation and Conventions}
\label{sec:notation}

\emph{Critical line.}
Throughout, the critical line means $\Re(s)=\tfrac12$.
The trace parameter $\sigma\in\mathbb{R}$ is evaluated at
$\sigma=\tfrac12$ unless stated otherwise.

\medskip
\emph{Global parameters.}
The reference values are $\kappa=53$ (prime cutoff, yielding
$16$ active primes $p\le\kappa$), $\varepsilon=0.05$ (Gaussian
regularisation parameter), and $N=100$ (number of Riemann
ordinates $\gamma_k$ employed).
These values are fixed throughout unless explicitly varied.
The key weight sums are
\begin{align*}
\Wgamma &\;:=\; \sum_{k=1}^N e^{-\varepsilon^2\gamma_k^2}\,\gamma_k,
\qquad
A(\varepsilon,N) \;:=\; \sum_{k=1}^N e^{-\varepsilon^2\gamma_k^2}, \\
P(\kappa) &\;:=\; \sum_{p\le\kappa}\log p
\qquad (\text{the Chebyshev theta sum}).
\end{align*}
At the reference parameters,
$\mathcal{W}_{0.05,100}\approx 28.42$ and $P(53)\approx 44.93$.

\medskip
\emph{Main objects.}
The coupling operator $\Phi(\sigma):\Hstr\to\Hnull$ is defined by
\[
\Phi(\sigma)_{k,p}=e^{-\varepsilon^2\gamma_k^2/2}\sin(\sigma\gamma_k\log p),
\]
and the Tehrani operator is $\Tt(\sigma)=\Phi(\sigma)\circ\Phi(\sigma)^*$.
The trace oscillatory sum is
\[
O(\sigma)=\tfrac12\sum_{k,p}e^{-\varepsilon^2\gamma_k^2}\cos(2\sigma\gamma_k\log p),
\]
with derivatives $O'(\sigma)$ and $O''(\sigma)$ taken with respect to $\sigma$.

\emph{Curvature and bias.}
The direct curvature sum is
\[
B:=\sum_{k,p}e^{-\varepsilon^2\gamma_k^2}(\gamma_k\log p)^2\cos(\gamma_k\log p),
\]
and the integrated bias is
\[
B_{\mathrm{int}}^+(\varepsilon,N):=\sum_{p\le\kappa}(\log p)^2 Z_{p,N}^+(\varepsilon).
\]

\emph{Ordinate proxies.}
The finite proxy used in numerics is
\[
Z_{p,N}^+(\varepsilon):=\mathrm{Re}\,Z_p(\varepsilon,N),
\qquad
Z_p(\varepsilon,N):=\sum_{k=1}^N e^{-\varepsilon^2\gamma_k^2}p^{i\gamma_k}.
\]
The asymptotic analogue is
\[
Z_{p,\infty}^+(\varepsilon):=\sum_{k\ge 1}e^{-\varepsilon^2\gamma_k^2}\cos(\gamma_k\log p),
\]
and the full signed asymptotic sum is
$Z_p^\infty(\varepsilon):=2Z_{p,\infty}^+(\varepsilon)$.
The full Guinand--Weil object, summed over all non-trivial zeros, is
\[
\widetilde Z_p^{\mathrm{GW,full}}(\varepsilon)
:=\sum_\rho h^{\mathrm{even}}_{p,\varepsilon}((\rho-\tfrac12)/i),
\]
where $h^{\mathrm{even}}_{p,\varepsilon}(u)=e^{-\varepsilon^2 u^2}\cos(u\log p)$.

\emph{Leading term and exponential sum.}
The Guinand--Weil leading term is
\[
\mathrm{Main}_p(\varepsilon):=
-\frac{\log p}{2\sqrt{\pi}\,\varepsilon\sqrt{p}}.
\]
The constant term $\mathrm{Const}_p^\infty$ used in
Theorem~\ref{thm:rpinf} is defined locally there.
The exponential sum over primes is
\[
M_k(\kappa):=\sum_{p\le\kappa}(\log p)\sin(\gamma_k\log p).
\]

\medskip
\emph{Analytical foundations.}
The proof of Theorem~\ref{thm:rpinf} uses the exact Guinand--Weil
decomposition of~\cite{Paper6} and classical Gaussian Fourier analysis.
Gaussian decay together with standard Riemann--von Mangoldt zero-counting
estimates~\cite[Thm.\,5.8]{IwaniecKowalski04} justifies convergence of
the infinite ordinate proxies and the $N\to\infty$ form of the
derivative bound; these estimates are not used as sign input in the
asymptotic proof of Theorem~\ref{thm:rpinf}.
Lemma~\ref{lem:subleading} is a separate prime-power Gaussian tail
estimate that does not invoke zero-density.
The conditional selection criterion (Proposition~\ref{prop:deriv},
Corollaries~\ref{cor:intbias_combined} and~\ref{cor:selection})
uses the trace formula of~\cite{Paper5}
and the curvature identity of~\cite{Paper6}.

\medskip
\emph{Numerical computations.}
All numerical results were computed with \texttt{mpmath} at 50~decimal
digits of working precision; Riemann zeros $\gamma_k$ were evaluated
via \texttt{mpmath.zetazero}.
Verification code is available at~\cite{Code}
(\url{https://github.com/utehrani/analysislab-nt}).

\medskip
\emph{Critical distinctions.}
The following symbols look similar but are distinct objects
(referenced throughout the paper):
\begin{center}
\renewcommand{\arraystretch}{1.25}
\begin{tabular}{l|l}
\hline
\textbf{NOT} & \textbf{BUT} \\
\hline
$Z_{p,N}^+$ (finite, numerics) & $Z_{p,\infty}^+$ (infinite, Thm~\ref{thm:rpinf}) \\
$Z_p^\infty=2Z_{p,\infty}^+$ (proxy, factor 2) & $\widetilde Z_p^{\mathrm{GW,full}}$ (full GW, all zeros) \\
$B$ ($\gamma_k^2$-weighted, $-19\,342.5$) & $B_{\mathrm{int}}^+$ (unweighted, $-42.21$) \\
$\mathcal{W}_{\varepsilon,N}$ ($\gamma_k$-weighted, this paper) & $A(\varepsilon,N)$ (unweighted, Paper~5) \\
$\mathrm{Const}_p$ (finite-$N$, Papers~5/6) & $\mathrm{Const}_p^\infty$ (this paper, Thm~\ref{thm:rpinf}) \\
$r_p(\varepsilon)=0.75$--$0.80$ at $\varepsilon=0.05$ & $r_p^\infty=\tfrac12$ ($\varepsilon\to 0^+$ limit; Thm~\ref{thm:rpinf}) \\
$D_{\mathrm{SEL}}\approx 10.985$ (Paper~5) & $D\approx 9.470$ (Paper~2) \\
\hline
\end{tabular}
\end{center}
Truncation: $|Z_{p,\infty}^+-Z_{p,N}^+|/|\mathrm{Main}_p|
\lesssim 10^{-61}$ at reference parameters
(truncation bound, Paper~6, \S\,3).
The identity $\widetilde Z_p^{\mathrm{GW,full}}=Z_p^\infty$ holds when all
non-trivial zeros lie on the critical line
(Paper~5, \S\,6; Assumption~\ref{ass:proxy}).

\vspace{2em}
\newpage

% ============================================================
\section{Introduction}
\label{sec:intro}
% ============================================================

\subsection{Background and Motivation}

This paper is the seventh in a series developing an operator-theoretic
framework for studying the distribution of the non-trivial zeros of
the Riemann zeta function.
The central construction is a finite-dimensional Hilbert-space model
in which prime numbers mediate coupling between zeros, encoded in an
explicit family of operators whose spectral properties reflect
arithmetic structure at the critical line.

\textbf{Paper~1}~\cite{Paper1} introduced the local curvature function
$\Hlocal(\sigma,\kappa)=\psi_1(\sigma)+\sum_{p\le\kappa}V_p(\sigma)$,
where $V_p(\sigma)$ is the prime-$p$ contribution to the spectral curvature.
At $\sigma=\tfrac12$ the sum diverges like $2(\log\kappa)^2$ as $\kappa\to\infty$,
identifying the critical line as a phase boundary through a structural
divergence of the curvature.

\textbf{Paper~2}~\cite{Paper2} connected this local curvature to the
Weil explicit functional, constructing $\sigma$-dependent admissible
Gaussian test functions $g^*_{\sigma,\varepsilon}$ for which the
curvature divergence translates into the growth of the diagonal
energy partial sums. In the renormalised critical-line scale,
$\sum_{p\le\kappa}f_p$ converges to the diagonal energy
$D:=\sum_p f_p\approx 9.470$.
The Weil positivity framework of~\cite{Lagarias} was identified as the
natural target for future operator-theoretic results.

\textbf{Paper~3}~\cite{Paper3} made the geometric structure explicit
by introducing two finite-dimensional Hilbert spaces, $\Hstr$ over
primes and $\Hnull$ over zeros, together with the coupling operator
$\Phi:\Hstr\to\Hnull$ defined by the prime-zero phasor expression
$\Phi_{k,p}=e^{-\varepsilon^2\gamma_k^2/2}\sin(\gamma_k\log p)$.
The later one-parameter convention $\sin(\sigma\gamma_k\log p)$
is introduced in Paper~5 and is the convention used throughout
the present paper.
The loop operator $T=\Phi^*\circ\Phi$ on $\Hstr$ and the energy
asymmetry ratio $\etaorig>0$ were introduced, with the latter
confirmed numerically.

\textbf{Paper~4}~\cite{Paper4} introduced the dual operator
$\Tt=\Phi\circ\Phi^*$ on $\Hnull$ --- the Tehrani operator --- and
proved conditional energy asymmetry under an explicit remainder-control
assumption.
A Hilbert--P\'olya diagnostic yielded $r_2\to 0.16$ as $\kappa\to\infty$
[numerical], numerically falsifying a Hilbert--P\'olya interpretation
of $\Tt$ on the tested grid and
distinguishing it from the constructions of Berry and
Keating~\cite{BerryKeating} and de Branges~\cite{deBranges}.

\textbf{Paper~5}~\cite{Paper5} extended the operator family to a
one-parameter family $\Tt(\sigma)=\Phi(\sigma)\circ\Phi(\sigma)^*$
and proved the exact spectral trace identity
\[
\mathrm{Tr}(\Tt(\sigma)) \;=\; \DSEL - O(\sigma), \qquad
O(\sigma) \;=\; \tfrac12\sum_{k,p}e^{-\varepsilon^2\gamma_k^2}
             \cos(2\sigma\gamma_k\log p),
\]
where $\DSEL$ is a $\sigma$-independent constant.
The paper introduced the Guinand--Weil smoothed zero sum
$\widetilde Z_p^{\mathrm{GW,full}}$ as formal motivation for analysing the
bias of prime--zero sums, and formulated the Bias Conjecture for the
infinite ordinate proxy: $Z_p^\infty(\varepsilon)<0$ for all primes
$p\le\kappa$ at small enough $\varepsilon$.
The transfer from the full Guinand--Weil object to the ordinate proxy
was left open as the GW bridge.

\textbf{Paper~6}~\cite{Paper6} proved the curvature--bias identity
$O''(\tfrac12)=-2B$ by direct differentiation of the trace formula of
Paper~5, where
$B=\sum_{k,p}e^{-\varepsilon^2\gamma_k^2}(\gamma_k\log p)^2\cos(\gamma_k\log p)$.
The leading term $\mathrm{Main}_p(\varepsilon)$ of the Guinand--Weil
decomposition was proved strictly negative for every prime $p$ and
every $\varepsilon>0$, and the numerical value $B=-19\,342.5$ at the
reference parameters gives $O''(\tfrac12)=+38\,685.1>0$.
Five open problems were formulated, of which the present paper addresses
the first two open problems of Paper~6:
conditional stationarity of $O'(\tfrac12)$ and the pointwise asymptotic bias of $Z_{p,\infty}^+$.

\smallskip
\noindent\emph{What is new in the present paper.}
Papers~5--6 \emph{conjectured} that $Z_p^\infty(\varepsilon)<0$ in
the limit $\varepsilon\to 0^+$ and tested it numerically.
The present paper \emph{proves the exact asymptotic rate}
$r_p^\infty=\tfrac12$ as a pointwise identity in $p$, conditional on
two named hypotheses (Assumptions~\ref{ass:subleading}
and~\ref{ass:proxy}), with an explicit decomposition of
$\mathrm{Const}_p^\infty$ into a pole term, the dominant
$-\tfrac12\mathrm{Main}_p$, an exponentially small reflected peak,
and a residual that is $o(|\mathrm{Main}_p|)$.
The contribution is the precise classification of orders of magnitude
in this identity --- not a new arithmetic conjecture.

The operator $\Tt(\sigma)$ constructed in this series is a Gram-type
operator with explicit arithmetic matrix entries; it is not a
Hilbert--P\'olya operator (Paper~4), and does not arise from an
adelic or non-commutative geometric construction.
Connes and Consani~\cite{ConnesConsani23} and Connes, Consani, and
Moscovici~\cite{CCM25} have constructed spectral triples and zeta-cycles
on adelic spaces whose spectra encode the zeros of $\zeta$ via abstract
operator-algebraic data.
Our framework is finite-cutoff, explicit, and operates within classical
Hilbert space theory over $\mathbb{C}^P$ and $\mathbb{C}^N$;
the bridge to the Weil positivity criterion of Lagarias~\cite{Lagarias}
and Bombieri~\cite{Bombieri} constitutes our principal open problem
(Open Problem~\ref{op:weil}).
Related approaches via the Hilbert space structure of the Weil
distribution have been pursued by Burnol~\cite{Burnol},
Suzuki~\cite{Suzuki25}, and others.

The present paper addresses the analytic question at the core of
Paper~6: whether a structural bound on $|O'(\tfrac12)|$
relative to $\Wgamma\cdot P(\kappa)$ can be derived from a
structural property of the prime-zero interaction.
Under the stated transfer and subleading hypotheses, we obtain
an asymptotic answer: the ratio
$r_p^\infty=\tfrac12$ is universal (Theorem~\ref{thm:rpinf}), and
the integrated-bias window is conditional on a subleading
archimedean estimate (Assumption~\ref{ass:subleading})
and a proxy-transfer condition (Assumption~\ref{ass:proxy}),
via Proposition~\ref{prop:intbias}.
The stationarity bound requires a logarithmic-saving hypothesis on
exponential sums, which we formulate precisely as Assumption~\Astar.
We have not obtained an unconditional bound of the form required by
Proposition~\ref{prop:deriv} from the five standard routes discussed
in \S\,\ref{sec:cancel}.
These are methodological observations on the present formulation,
not impossibility results.
Assumption~\Astar\ is therefore taken as the natural conditional input
(\S\,\ref{sec:cancel}, Remark~\ref{rem:necessity}).

\subsection{Main Results}

The four main results of this paper are as follows.

The first, in \S\,\ref{sec:rpinf} (Theorem~\ref{thm:rpinf}), is the
\emph{pointwise asymptotic limit}: for every prime $p\ge 2$,
$r_p^\infty := \lim_{\varepsilon\to 0^+}
|\mathrm{Const}_p^\infty(\varepsilon)|/|\mathrm{Main}_p(\varepsilon)| = \tfrac12$.
This is conditional on Assumptions~\ref{ass:subleading} and~\ref{ass:proxy},
and resolves the pointwise asymptotic bias question
of Paper~6, \S\,8, for the ordinate proxy; the GW bridge
to the full object $\widetilde Z_p^{\mathrm{GW,full}}$ remains open
(Open Problem~\ref{op:gwbridge}).

The second, in \S\,\ref{sec:decomp} (Definition~\ref{def:threeterm}), is
the \emph{three-term decomposition} of $O'(\tfrac12)$: the algebraic
definition $O'(\tfrac12)=T_{\mathrm{pol}}(\varepsilon)+T_{\mathrm{main}}(\varepsilon)+T_{\mathrm{res}}(\kappa,\varepsilon,N)$
where $T_{\mathrm{pol}}(\varepsilon)$ and $T_{\mathrm{main}}(\varepsilon)$ arise from the GW pole and
prime main-term contributions and are algebraically explicit, while $T_{\mathrm{res}}$ is the
residual (Definition~\ref{def:threeterm}).
At the reference parameters, $|O'(\tfrac12)|=2.4751$ is only
approximately $0.2\%$ of the trivial bound $\Wgamma\cdot P(\kappa)\approx 1277$,
with $T_{\mathrm{pol}}+T_{\mathrm{main}}=+38.84$ and $T_{\mathrm{res}}=-36.37$ cancelling to a factor of $15.7$
above $|O'(\tfrac12)|$ [numerical].

The third, in \S\,\ref{sec:curv} (Proposition~\ref{prop:intbias}), is
the \emph{integrated-bias window}: for every prime cutoff $\kappa\ge 2$
there exists $\varepsilon_{\mathrm{int}}(\kappa)>0$ such that
$B_{\mathrm{int}}^\infty(\kappa,\varepsilon)<0$ for all
$\varepsilon\in(0,\varepsilon_{\mathrm{int}}(\kappa))$.
This result is conditional on Assumptions~\ref{ass:subleading}
and~\ref{ass:proxy}, and addresses the curvature-uniformity question
of Paper~6, \S\,8.
At reference parameters $O''(\tfrac12)=+38\,685.1>0$ [numerical].

The fourth, in \S\,\ref{sec:sel} (Proposition~\ref{prop:deriv} and
Corollaries~\ref{cor:intbias_combined},~\ref{cor:selection}), is the
\emph{conditional selection criterion}: under Assumption~\Astar\
(Cancellation Hypothesis, \S\,\ref{sec:cancel}),
$|O'(\tfrac12;\kappa,\varepsilon,N)|\le C\cdot\Wgamma\cdot P(\kappa)$
where $C:=C_0/(\log\gamma_1)^A$
and $\log\gamma_1\approx 2.6486$ is independent of $\kappa$, $\varepsilon$, and~$N$.
For $\varepsilon\in(0,\varepsilon_{\mathrm{int}}(\kappa))$ this gives a
conditional derivative bound at $\sigma=\tfrac12$
(Proposition~\ref{prop:deriv}); genuine smallness
requires the additional condition $C_0/(\log\gamma_1)^A<1$.
We make no claim that this condition is automatically met:
verifying or proving $C_0/(\log\gamma_1)^A<1$ for an explicit
$(C_0,A)$ is itself part of the analytical content of
Assumption~\Astar.
The present paper isolates this requirement; resolving it is left
open (see~\S\,\ref{sec:open}).
Combined with $B_{\mathrm{int}}^\infty<0$
(Corollary~\ref{cor:intbias_combined}) and the additional hypotheses
$B(\kappa,\varepsilon,N)<0$ and exact stationarity $O'(\tfrac12)=0$,
this yields a conditional local-selection criterion at
$\sigma=\tfrac12$ (Corollary~\ref{cor:selection}).

\subsection{Reader Contract}

This paper \emph{establishes unconditionally}:
Definition~\ref{def:threeterm} (bookkeeping decomposition);
Proposition~\ref{prop:antisym} (antisymmetry).
Observation~\ref{obs:threeterm} (numerical values at reference parameters)
is numerical, not a proved statement.

This paper \emph{proves under
Assumptions~\ref{ass:subleading} (numerically supported, analytically open)
and~\ref{ass:proxy} (proxy-transfer, formally weaker in content
than RH: RH implies it, no converse claimed)}:
Theorem~\ref{thm:rpinf} ($r_p^\infty=\tfrac12$ for every prime $p$);
Corollary~\ref{cor:sign} ($Z_{p,\infty}^+(\varepsilon)<0$ for small $\varepsilon$);
Proposition~\ref{prop:intbias}
(integrated-bias window $\varepsilon_{\mathrm{int}}(\kappa)>0$ for all $\kappa\ge 2$).

Under Assumption~\Astar\
(\S\,\ref{sec:sel}, Proposition~\ref{prop:deriv}):
$|O'(\tfrac12)|\le C\cdot\Wgamma\cdot P(\kappa)$.
Under the additional hypotheses $B(\kappa,\varepsilon,N)<0$ and
$O'(\tfrac12)=0$ (Corollary~\ref{cor:selection}):
$\sigma=\tfrac12$ is a strict local minimum of $O$.

This paper \emph{establishes numerically} at
$(\kappa,\varepsilon,N)=(53,0.05,100)$:
$O'(\tfrac12)=+2.4751$; $O''(\tfrac12)=+38\,685.1$; $B=-19\,342.5$;
$\varepsilon_{\mathrm{int}}(53)>0.200$; cancellation factor $15.7\times$ below trivial bound.

This paper \emph{leaves open}:
the Weighted Bias Bridge (Open Problem~\ref{op:wbridge});
Assumption~\Astar\ (Open Problem~\ref{op:cancel});
the Guinand--Weil bridge with off-critical control (Open Problem~\ref{op:gwbridge});
the scaling of $\varepsilon_{\mathrm{int}}(\kappa)$ (Open Problem~\ref{op:epsc});
the near-minimum question (Open Problem~\ref{op:global});
and the Weil-transfer connection (Open Problem~\ref{op:weil}).

% ============================================================
\section{Background: Trace Formula and Guinand--Weil Framework}
\label{sec:bg}
% ============================================================

\subsection{The Operator Family and Trace Formula}

The two-parameter coupling operator $\Phi(\sigma):\Hstr\to\Hnull$
is defined on the finite-dimensional spaces $\Hstr\cong\mathbb{C}^P$
(over primes $p\le\kappa$) and $\Hnull\cong\mathbb{C}^N$
(over zeros $\gamma_1,\ldots,\gamma_N$) by
\[
\Phi(\sigma)_{k,p} \;=\; e^{-\varepsilon^2\gamma_k^2/2}\,\sin(\sigma\gamma_k\log p).
\]
The Tehrani operator is $\Tt(\sigma)=\Phi(\sigma)\circ\Phi(\sigma)^*$,
a self-adjoint endomorphism of $\Hnull$.
Paper~5~\cite{Paper5} proves the following exact identity.

\begin{theorem}[Trace formula; proved in \cite{Paper5}]
\label{thm:trace}
For all $\sigma\in\mathbb{R}$,
\[
\mathrm{Tr}(\Tt(\sigma)) \;=\; \DSEL - O(\sigma),
\]
where $\DSEL:=\tfrac12 A(\varepsilon,N)\cdot\pi(\kappa)$ is
$\sigma$-independent and
\[
O(\sigma):=\tfrac12\sum_{k=1}^N\sum_{p\le\kappa}
e^{-\varepsilon^2\gamma_k^2}\cos(2\sigma\gamma_k\log p).
\]
At the reference parameters $\DSEL\approx 10.985$.
\end{theorem}

Differentiating twice in $\sigma$ yields the curvature
\[
O''(\sigma)=-2\sum_{k,p}e^{-\varepsilon^2\gamma_k^2}
(\gamma_k\log p)^2\cos(2\sigma\gamma_k\log p),
\]
and at $\sigma=\tfrac12$ this equals $-2B$ by Paper~6~\cite{Paper6}.

\subsection{The Curvature Identity}

\begin{theorem}[Curvature--bias identity; proved in \cite{Paper6}]
\label{thm:curvbias}
The direct curvature sum $B:=\sum_{k,p}e^{-\varepsilon^2\gamma_k^2}(\gamma_k\log p)^2\cos(\gamma_k\log p)$
satisfies the algebraic identity
\[
O''\!\left(\tfrac12\right) \;=\; -2B.
\]
At the reference parameters $B=-19\,342.5$ and $O''(\tfrac12)=+38\,685.1$ [numerical].
\end{theorem}

\subsection{The Guinand--Weil Framework}

The Guinand--Weil explicit formula~\cite{Guinand,Weil} applied to the
even test function $h^{\mathrm{even}}_{p,\varepsilon}(u)=e^{-\varepsilon^2 u^2}\cos(u\log p)$
yields a decomposition of the formal sum over non-trivial zeros,
following the conventions of~\cite{IwaniecKowalski04}.
For the ordinate proxy $Z_{p,N}^+(\varepsilon)=\mathrm{Re}\,Z_p(\varepsilon,N)$
at finite truncation level~$N$ one obtains
\begin{equation}\label{eq:gw-decomp}
Z_{p,N}^+(\varepsilon)
  = \underbrace{\tfrac12\,\mathrm{POL}(\varepsilon,p)
    - \mathrm{PRIMS}_{n=p}(\varepsilon)}_{\text{leading}}
    - \mathrm{PRIMS}_{n\ne p}(\varepsilon)
    + \Gamma_p(\varepsilon)
    + \tfrac12\,\mathrm{CST}(\varepsilon)
    + R_{p,N}(\varepsilon).
\end{equation}
The leading terms are:
$\tfrac12\,\mathrm{POL}(\varepsilon,p)=e^{\varepsilon^2/4}\cosh(\log p/2)$
(Lemma~\ref{lem:pol});
$-\mathrm{PRIMS}_{n=p}(\varepsilon)
=\mathrm{Main}_p(\varepsilon)\cdot\tfrac12(1+e^{-(\log p)^2/\varepsilon^2})$
(Lemma~\ref{lem:main}),
where $\mathrm{Main}_p(\varepsilon):=-({\log p})/({2\sqrt{\pi}\,\varepsilon\sqrt{p}})$.
The off-diagonal contribution $\mathrm{PRIMS}_{n\ne p}$ is exponentially
small (Lemma~\ref{lem:subleading}).
The archimedean contribution $\Gamma_p$ and constant term
$\tfrac12\,\mathrm{CST}$ are subleading under
Assumption~\ref{ass:subleading}.
The remainder $R_{p,N}$ is negligible at reference parameters
(truncation estimate of~\cite{Paper6}).

For Theorem~\ref{thm:rpinf} we work with the infinite analogue
\[
\mathrm{Const}_p^\infty(\varepsilon)
  := Z_{p,\infty}^+(\varepsilon) - \mathrm{Main}_p(\varepsilon)
     - \Gamma_p(\varepsilon),
\]
where $Z_{p,\infty}^+(\varepsilon):=\sum_{k\ge 1}
e^{-\varepsilon^2\gamma_k^2}\cos(\gamma_k\log p)$.

\begin{remark}
The identity $\widetilde Z_p^{\mathrm{GW,full}}(\varepsilon)=Z_p^\infty(\varepsilon)$
holds when all non-trivial zeros lie on the critical line; without
this assumption an off-critical contribution must be added (Paper~5, \S\,6).
All computations in this paper use the ordinate proxy $Z_{p,N}^+(\varepsilon)$.
\end{remark}

% ============================================================
\section{The Pointwise Asymptotic Limit}
\label{sec:rpinf}
% ============================================================

\begin{theorem}[Pointwise asymptotic limit; proved under
                 Assumptions~\emph{\ref{ass:subleading}} and~\emph{\ref{ass:proxy}}]
\label{thm:rpinf}
For every prime $p\ge 2$,
\[
r_p^\infty \;:=\; \lim_{\varepsilon\to 0^+}
  \frac{|\mathrm{Const}_p^\infty(\varepsilon)|}
       {|\mathrm{Main}_p(\varepsilon)|}
  \;=\; \frac{1}{2},
\]
where $\mathrm{Const}_p^\infty(\varepsilon):=Z_{p,\infty}^+(\varepsilon)
-\mathrm{Main}_p(\varepsilon)-\Gamma_p(\varepsilon)$.
The limit is independent of $p$.
\end{theorem}

\begin{assumption}[Subleading archimedean terms]
\label{ass:subleading}
For every prime $p\ge 2$, as $\varepsilon\to0^+$,
\[
  |\Gamma_p(\varepsilon)|
  +\tfrac12|\mathrm{CST}(\varepsilon)|
  = o\!\left(|\mathrm{Main}_p(\varepsilon)|\right).
\]
\end{assumption}

\begin{remark}
Assumption~\ref{ass:subleading} is numerically supported but not
proved. For the gamma contribution, Paper~6 reports finite-grid
numerical evidence that the ratio
$|\Gamma_p(\varepsilon)|/|\mathrm{Main}_p(\varepsilon)|$
is consistent with linear scaling in~$\varepsilon$ on the tested
grid: across all $16$ primes $p\le 53$ and five values
$\varepsilon\in\{0.02,0.03,0.05,0.08,0.12\}$, the log-log slope
is $1.001\pm 0.003$.
This is evidence for
$|\Gamma_p(\varepsilon)|/|\mathrm{Main}_p(\varepsilon)|
=O_{\mathrm{num}}(\varepsilon)$,
not an analytic proof of a limit as $\varepsilon\to 0^+$.
Since $|\mathrm{Main}_p(\varepsilon)|\asymp\varepsilon^{-1}$,
ratio scaling of this form would correspond to
$|\Gamma_p(\varepsilon)|=O_{\mathrm{num}}(1)$,
not to $|\Gamma_p(\varepsilon)|=O(\varepsilon)$.
The estimate
$|\Gamma_p(\varepsilon)|+\tfrac12|\mathrm{CST}(\varepsilon)|
=o(|\mathrm{Main}_p(\varepsilon)|)$
therefore remains exactly the hypothesis of
Assumption~\ref{ass:subleading}.
The $o$ statement for the constant-term contribution
$|\mathrm{CST}(\varepsilon)|/|\mathrm{Main}_p(\varepsilon)|\to 0$
is analytically open; numerically
$|\mathrm{CST}(\varepsilon)|/|\mathrm{Main}_p(\varepsilon)|\sim\varepsilon$
[numerical, same dataset as $\Gamma_p$].
An analytic proof of the combined estimate remains open.
\end{remark}

\begin{assumption}[Proxy transfer]
\label{ass:proxy}
For every fixed prime $p\ge 2$, the full Guinand--Weil zero
sum admits a decomposition
\[
\widetilde Z_p^{\mathrm{GW,full}}(\varepsilon)
  = Z_p^\infty(\varepsilon) + E_p^{\mathrm{off}}(\varepsilon),
\qquad
Z_p^\infty(\varepsilon)
  := 2 Z_{p,\infty}^+(\varepsilon),
\]
where
\[
E_p^{\mathrm{off}}(\varepsilon)
  = o(|\mathrm{Main}_p(\varepsilon)|)
  \qquad (\varepsilon\to 0^+).
\]
Equivalently, the positive-ordinate proxy satisfies
\[
Z_{p,\infty}^+(\varepsilon)
  = \tfrac12\,\widetilde Z_p^{\mathrm{GW,full}}(\varepsilon)
    - \tfrac12\,E_p^{\mathrm{off}}(\varepsilon).
\]
This half-Guinand--Weil form is used in
Theorem~\ref{thm:rpinf}.
\end{assumption}

\begin{remark}
Assumption~\ref{ass:proxy} is implied by RH and is formally
weaker in content: it requires only asymptotic smallness of the
off-critical contribution, not its vanishing.
No converse implication is claimed.
No analytic proof is available without RH.
\end{remark}

The proof rests on three lemmas which identify the leading terms of
$\mathrm{Const}_p^\infty(\varepsilon)$ as $\varepsilon\to 0$.

\begin{lemma}[Pole term; proved]
\label{lem:pol}
For every prime $p\ge 2$ and every $\varepsilon>0$,
\[
\tfrac12\,\mathrm{POL}(\varepsilon,p)
  \;=\; e^{\varepsilon^2/4}\cosh\!\left(\tfrac{\log p}{2}\right).
\]
\end{lemma}

\begin{proof}
By definition,
$h^{\mathrm{even}}_{p,\varepsilon}(i/2)=e^{-\varepsilon^2(i/2)^2}\cos(i\log p/2)
=e^{\varepsilon^2/4}\cosh(\log p/2)$.
The evaluation at $u=-i/2$ gives the same value by the even symmetry
of $h^{\mathrm{even}}_{p,\varepsilon}$.
The pole contribution in the Guinand--Weil formula is
$(h(i/2)+h(-i/2))/2=e^{\varepsilon^2/4}\cosh(\log p/2)$.
\end{proof}

\begin{lemma}[Prime-$p$ main term; proved]
\label{lem:main}
For every prime $p\ge 2$ and every $\varepsilon>0$,
\[
-\mathrm{PRIMS}_{n=p}(\varepsilon)
  \;=\; \tfrac12\,\mathrm{Main}_p(\varepsilon)\cdot
        \bigl[1 + e^{-(\log p)^2/\varepsilon^2}\bigr],
\]
where $\mathrm{PRIMS}_{n=p}(\varepsilon)$ denotes the $n=p$ term
in the prime-power sum of the Guinand--Weil formula.
\end{lemma}

\begin{proof}
For the positive-ordinate proxy $Z_{p,\infty}^+$, the
Guinand--Weil prime term at $n=p$ in the convention
of~\cite[Ch.\,5]{IwaniecKowalski04} contributes
\[
-\frac{1}{2\pi}\,\frac{\Lambda(p)}{\sqrt p}\,
\hat h^{\mathrm{even}}_{p,\varepsilon}
\!\left(\frac{\log p}{2\pi}\right).
\]
Substituting the Fourier evaluation
$\hat h^{\mathrm{even}}_{p,\varepsilon}(\log p/(2\pi))
=\frac{\sqrt\pi}{2\varepsilon}\bigl(1+e^{-(\log p)^2/\varepsilon^2}\bigr)$
and $\Lambda(p)=\log p$:
\[
-\frac{1}{2\pi}\cdot\frac{\log p}{\sqrt p}\cdot
\frac{\sqrt\pi}{2\varepsilon}
\bigl(1+e^{-(\log p)^2/\varepsilon^2}\bigr)
= -\frac{\log p}{4\sqrt\pi\,\varepsilon\sqrt p}
\bigl(1+e^{-(\log p)^2/\varepsilon^2}\bigr)
= \tfrac12\,\mathrm{Main}_p(\varepsilon)
\bigl(1+e^{-(\log p)^2/\varepsilon^2}\bigr),
\]
where the last equality uses
$\mathrm{Main}_p(\varepsilon)
=-(\log p)/(2\sqrt\pi\,\varepsilon\sqrt p)$.
\end{proof}

\begin{lemma}[Off-diagonal prime-power bound; proved]
\label{lem:subleading}
For the off-diagonal prime-power contributions there exist
constants $C(p)>0$ and $c_p>0$, depending only on $p$,
such that for all $\varepsilon\in(0,1)$,
\[
|\mathrm{PRIMS}_{n\ne p}(\varepsilon)|
  \;\le\; \frac{C(p)}{\varepsilon}\,e^{-c_p/\varepsilon^2},
\]
where $c_p:=\min(d_p^2,(2\log 2)^2)/4>0$ and
$d_p:=\min\{|\log(n/p)|:n=q^k,\,n\ne p\}>0$.
\end{lemma}

\begin{remark}
The archimedean and constant-term subleading estimate
$|\Gamma_p(\varepsilon)|+\tfrac12|\mathrm{CST}(\varepsilon)|
=o(|\mathrm{Main}_p(\varepsilon)|)$
is the content of Assumption~\ref{ass:subleading};
see the numerical verification in the remark following that assumption.
\end{remark}

\begin{proof}
Write $a:=\log p$.
The prime-power sum $\mathrm{PRIMS}_{n\ne p}$ is bounded by
$\sum_{n\ne p}\Lambda(n)/\sqrt{n}\cdot|\hat h_{p,\varepsilon}(\log n/(2\pi))|$,
and the Fourier transform of $h_{p,\varepsilon}^{\mathrm{even}}$ satisfies
\[
|\hat h_{p,\varepsilon}(\log n/(2\pi))|
\le \frac{\sqrt\pi}{2\varepsilon}\!
\left[e^{-(\log n-a)^2/(4\varepsilon^2)}
+e^{-(\log n+a)^2/(4\varepsilon^2)}\right].
\]
Split the prime-power sum into the near range
$\mathcal{N}_p:=\{n=q^m:n\ne p,\,|\log n-a|\le 1\}$
and its complement.

\emph{Near range.}
$\mathcal{N}_p$ is finite; for each $n\in\mathcal{N}_p$,
$|\log n-a|\ge d_p>0$, so its contribution is
$O_p(\varepsilon^{-1}e^{-d_p^2/(4\varepsilon^2)})$.

\emph{Far range $|\log n - a|>1$.}
For $\varepsilon<1$,
\[
e^{-(\log n-a)^2/(4\varepsilon^2)}
\le e^{-1/(8\varepsilon^2)}\,e^{-(\log n-a)^2/8}.
\]
The series $\sum_{n=q^m}\Lambda(n)/\sqrt{n}\cdot e^{-(\log n-a)^2/8}$
converges, since the Gaussian factor $e^{-(\log n)^2/8}$
decays faster than every power of~$n$.
Hence the far-range contribution is
$O_p(\varepsilon^{-1}e^{-1/(8\varepsilon^2)})$.

\emph{Reflected peak.}
The term $e^{-(\log n+a)^2/(4\varepsilon^2)}$ satisfies
$\log n+a\ge 2\log 2$ for all $n,p\ge 2$, so the reflected
contribution is bounded by
$e^{-(2\log 2)^2/(4\varepsilon^2)}$
times the same summable majorant.

Combining the three pieces and decreasing the exponent
constant if necessary gives the bound with
$c_p:=\min(d_p^2,(2\log 2)^2)/4>0$.
\end{proof}

For use in the proof of Theorem~\ref{thm:rpinf}, we record the
infinite half-sum form of the Guinand--Weil decomposition.
Combining the full Guinand--Weil formula~\eqref{eq:gw-decomp}
with Assumption~\ref{ass:proxy} gives, for fixed prime $p\ge 2$,
\[
Z_{p,\infty}^+(\varepsilon)
 \;=\; \tfrac12\,\mathrm{POL}(\varepsilon,p)
       \;-\; \mathrm{PRIMS}_{n=p}(\varepsilon)
       \;-\; \mathrm{PRIMS}_{n\ne p}(\varepsilon)
       \;+\; \Gamma_p(\varepsilon)
       \;+\; \tfrac12\,\mathrm{CST}(\varepsilon)
       \;-\; \tfrac12\,E_p^{\mathrm{off}}(\varepsilon).
\]
No finite-$N$ truncation term appears; convergence of the
positive-ordinate Gaussian sum is justified by Gaussian decay
together with the standard Riemann--von Mangoldt zero-counting
estimate (see \S\,\ref{sec:notation}, \emph{Analytical foundations}).

\smallskip
\noindent\emph{Strategy of the proof.}
The argument proceeds in three steps.
\emph{Step 1} substitutes the infinite half-sum form into the
definition of $\mathrm{Const}_p^\infty$, yielding an exact
identity with a pole term, a leading $-\tfrac12\mathrm{Main}_p$
contribution, an exponentially small reflected peak, and a
residual $R'$.
\emph{Step 2} bounds the residual $|R'(p,\varepsilon)|
=o(|\mathrm{Main}_p|)$ using Lemma~\ref{lem:subleading}
and Assumptions~\ref{ass:subleading},~\ref{ass:proxy}.
\emph{Step 3} computes the limit $\varepsilon\to 0^+$:
the pole term is $O(1)$ against $|\mathrm{Main}_p|\sim\varepsilon^{-1}$,
the reflected peak vanishes, and the residual is subleading by Step~2.
The dominant term is $-\tfrac12\mathrm{Main}_p$, giving
$r_p^\infty=\tfrac12$.

\begin{proof}[Proof of Theorem~\ref{thm:rpinf}]
\textit{Step 1: Exact identity.}
We work with the infinite proxy $Z_{p,\infty}^+$ and
$\mathrm{Const}_p^\infty(\varepsilon)
=Z_{p,\infty}^+(\varepsilon)-\mathrm{Main}_p(\varepsilon)-\Gamma_p(\varepsilon)$.
Applying the Guinand--Weil decomposition~\eqref{eq:gw-decomp} to
$\widetilde Z_p^{\mathrm{GW,full}}$ and then using
Assumption~\ref{ass:proxy} to transfer to the half-sum
$Z_{p,\infty}^+$ (up to an $o(|\mathrm{Main}_p|)$ error), as
recorded in the preamble above, together with
Lemmas~\ref{lem:pol} and~\ref{lem:main}:
\[
\mathrm{Const}_p^\infty(\varepsilon)
  \;=\; \underbrace{e^{\varepsilon^2/4}\cosh\!\left(\tfrac{\log p}{2}\right)}_{\text{pole, } O(1)}
        \;\underbrace{-\, \tfrac12\,\mathrm{Main}_p(\varepsilon)}_{\text{leading, } \sim\varepsilon^{-1}}
        \;+\; \underbrace{\tfrac12\,\mathrm{Main}_p(\varepsilon)\,
              e^{-(\log p)^2/\varepsilon^2}}_{\text{reflected peak, } \to 0}
        \;+\; \underbrace{R'(p,\varepsilon)}_{=\,o(|\mathrm{Main}_p|)},
\]
where $R'(p,\varepsilon)$ collects the off-diagonal prime-power term
$-\mathrm{PRIMS}_{n\ne p}(\varepsilon)$, the constant-term contribution
$\tfrac12\,\mathrm{CST}(\varepsilon)$, and the proxy-transfer error
$-\tfrac12\,E_p^{\mathrm{off}}(\varepsilon)$ from
Assumption~\ref{ass:proxy}.
(The $\Gamma_p(\varepsilon)$ contribution from the
Guinand--Weil decomposition is absorbed into the definition of
$\mathrm{Const}_p^\infty$ and does not appear in $R'$.)

\textit{Step 2: Bounding the residual.}
By Lemma~\ref{lem:subleading},
$|\mathrm{PRIMS}_{n\ne p}|\le C(p)\varepsilon^{-1}e^{-c_p/\varepsilon^2}$;
by Assumption~\ref{ass:subleading},
$\tfrac12|\mathrm{CST}(\varepsilon)|
\le|\Gamma_p(\varepsilon)|+\tfrac12|\mathrm{CST}(\varepsilon)|
=o(|\mathrm{Main}_p(\varepsilon)|)$;
by Assumption~\ref{ass:proxy},
$|E_p^{\mathrm{off}}(\varepsilon)|=o(|\mathrm{Main}_p(\varepsilon)|)$.
Together these give $|R'(p,\varepsilon)|=o(|\mathrm{Main}_p(\varepsilon)|)$
as $\varepsilon\to 0^+$.

\textit{Step 3: Taking the limit.}
As $\varepsilon\to 0^+$: $e^{\varepsilon^2/4}\to 1$, $e^{-(\log p)^2/\varepsilon^2}\to 0$, so
\[
\mathrm{Const}_p^\infty(\varepsilon)
  \;=\; \cosh\!\left(\tfrac{\log p}{2}\right)
        - \tfrac12\,\mathrm{Main}_p(\varepsilon)
        + o\!\left(|\mathrm{Main}_p(\varepsilon)|\right).
\]
Dividing by $|\mathrm{Main}_p(\varepsilon)| = (\log p)/(2\sqrt{\pi}\,\varepsilon\sqrt{p})\to\infty$:
\[
\frac{\mathrm{Const}_p^\infty(\varepsilon)}{\mathrm{Main}_p(\varepsilon)}
  \;\longrightarrow\;
  \frac{\cosh(\log p/2)}{\mathrm{Main}_p(\varepsilon)} - \tfrac12
  \;\longrightarrow\; 0 - \tfrac12 \;=\; -\tfrac12.
\]
Therefore $|\mathrm{Const}_p^\infty(\varepsilon)|/|\mathrm{Main}_p(\varepsilon)|
\to\tfrac12$ as $\varepsilon\to 0^+$, hence $r_p^\infty=\tfrac12$.
\end{proof}

\begin{corollary}[Sign of asymptotic proxy; proved under
                  Assumptions~\ref{ass:subleading} and~\ref{ass:proxy}]
\label{cor:sign}
For every prime $p\ge 2$ there exists $\varepsilon_0(p)>0$ such that
$Z_{p,\infty}^+(\varepsilon)<0$ for all $\varepsilon\in(0,\varepsilon_0(p))$.
\end{corollary}

\begin{proof}
By Theorem~\ref{thm:rpinf} and Assumption~\ref{ass:subleading},
\[
\frac{|\mathrm{Const}_p^\infty(\varepsilon)|}{|\mathrm{Main}_p(\varepsilon)|}
  \to\tfrac12,
\qquad
\frac{|\Gamma_p(\varepsilon)|}{|\mathrm{Main}_p(\varepsilon)|}
  \to 0
\qquad(\varepsilon\to 0^+).
\]
Choose $\varepsilon>0$ small enough that
\[
|\mathrm{Const}_p^\infty(\varepsilon)|
  \le \tfrac34 |\mathrm{Main}_p(\varepsilon)|,
\qquad
|\Gamma_p(\varepsilon)|
  \le \tfrac18 |\mathrm{Main}_p(\varepsilon)|.
\]
Since $\mathrm{Main}_p(\varepsilon)<0$, it follows that
\[
Z_{p,\infty}^+(\varepsilon)
=\mathrm{Main}_p(\varepsilon)+\Gamma_p(\varepsilon)
 +\mathrm{Const}_p^\infty(\varepsilon)
\le -\bigl(1-\tfrac34-\tfrac18\bigr)|\mathrm{Main}_p(\varepsilon)|<0.
\]
\end{proof}

% ============================================================
\section{\texorpdfstring{Three-Term Decomposition of $O'(\tfrac12)$}{Three-Term Decomposition of O'(1/2)}}
\label{sec:decomp}
% ============================================================

\begin{definition}[Three-term decomposition of $O'(\tfrac12)$]
\label{def:threeterm}
Define
\begin{align*}
T_{\mathrm{pol}}(\varepsilon) &\;:=\;
  \tfrac14\sum_{p\le\kappa}(\log p)\,e^{\varepsilon^2/4}
  \sinh\!\left(\tfrac{\log p}{2}\right), \\[4pt]
T_{\mathrm{main}}(\varepsilon) &\;:=\;
  -\sum_{p\le\kappa}(\log p)\,
  \frac{1-\tfrac12\log p}{4\sqrt{\pi}\,\varepsilon\sqrt{p}}, \\[4pt]
T_{\mathrm{res}}(\kappa,\varepsilon,N) &\;:=\;
  O'\!\left(\tfrac12;\kappa,\varepsilon,N\right)
  - T_{\mathrm{pol}}(\varepsilon) - T_{\mathrm{main}}(\varepsilon).
\end{align*}
The terms $T_{\mathrm{pol}}(\varepsilon)$ and $T_{\mathrm{main}}(\varepsilon)$ are
algebraically explicit in $\varepsilon$ and arise from the Guinand--Weil
pole and prime main-term contributions, respectively
(via Lemma~\ref{lem:pol} and the derivative
$\tfrac{\mathrm{d}}{\mathrm{d}\log p}[(\log p)/\sqrt{p}]
=(1-\tfrac12\log p)/\sqrt{p}$).
The residual $T_{\mathrm{res}}$ collects all remaining terms.
\end{definition}

\begin{observation}[Numerical values; numerical]
\label{obs:threeterm}
At $(\kappa,\varepsilon,N)=(53,0.05,100)$:
$T_{\mathrm{pol}}(\varepsilon)=+27.67$,
$T_{\mathrm{main}}(\varepsilon)=+11.18$,
$T_{\mathrm{res}}(53,0.05,100)=-36.37$,
with sum $O'(\tfrac12)=+2.4751$.
The trivial bound
$|O'(\tfrac12)|\le\Wgamma\cdot P(\kappa)\approx 28.42\times 44.93\approx 1277$
exceeds $|O'(\tfrac12)|$ by a factor of approximately $516$.
The algebraic terms satisfy $T_{\mathrm{pol}}+T_{\mathrm{main}}=+38.84$,
while $T_{\mathrm{res}}=-36.37$;
their partial cancellation leaves $|O'(\tfrac12)|=2.4751$, which is
a factor of $38.84/2.4751\approx 15.7$ below $|T_{\mathrm{pol}}+T_{\mathrm{main}}|$ [numerical].
\end{observation}

\begin{remark}[Domain]
Definition~\ref{def:threeterm} uses the ordinate proxy $Z_{p,N}^+$,
not the full Guinand--Weil object $\widetilde Z_p^{\mathrm{GW,full}}$;
the GW bridge remains open (Open Problem~\ref{op:gwbridge}).
\end{remark}

% ============================================================
\section{The Cancellation Hypothesis}
\label{sec:cancel}
% ============================================================

The terms $T_{\mathrm{pol}}(\varepsilon)$ and $T_{\mathrm{main}}(\varepsilon)$
of Definition~\ref{def:threeterm}
are explicitly computable and grow with $\kappa$ (they are sums over
all primes $p\le\kappa$).
The residual $T_{\mathrm{res}}(\kappa,\varepsilon,N)$ depends on the fine arithmetic
structure of the prime-zero interactions through the exponential sum
$M_k(\kappa)=\sum_{p\le\kappa}(\log p)\sin(\gamma_k\log p)$.
The following hypothesis asserts a logarithmic saving in $M_k(\kappa)$
relative to its trivial bound $P(\kappa)$.

\begin{assumption}[Cancellation Hypothesis \Astar]
\label{ass:cancel}
There exist absolute constants $C_0>0$ and $A\ge 1$ such that
for every Riemann ordinate $\gamma\in\{\gamma_k:k\ge 1\}$ and every
prime cutoff $\kappa\ge 2$,
\[
\left|\sum_{p\le\kappa}(\log p)\sin(\gamma\log p)\right|
  \;\le\; \frac{C_0\cdot P(\kappa)}{(\log\gamma)^A}.
\]
\end{assumption}

\smallskip
\noindent\emph{Two perspectives on $(A^{**})$.}
The hypothesis $|M_k(\kappa)|\le C\cdot P(\kappa)/(\log\gamma_k)^A$
admits two equivalent readings of the same sum
$M_k(\kappa)=\sum_{p\le\kappa}(\log p)\sin(\gamma_k\log p)$:

\smallskip
\noindent\emph{Primes interfere at a fixed zero.}
For each $\gamma_k$, the contributions $(\log p)\sin(\gamma_k\log p)$
from individual primes have varying signs.
$(A^{**})$ asserts that these contributions cancel beyond the
trivial $P(\kappa)$ envelope by a logarithmic margin.

\smallskip
\noindent\emph{A zero scans the prime structure.}
Writing $M_k(\kappa)=\mathrm{Im}\,\sum_{p\le\kappa}(\log p)\,p^{i\gamma_k}$,
each prime acquires a phase via the ordinate $\gamma_k$.
$(A^{**})$ asserts that no zero frequency creates strong coherent
alignment among the prime phases.

\smallskip
Both readings express the same arithmetic content:
\emph{no zero frequency aligns the primes coherently}.
This is the analytical question that the present paper isolates
but does not resolve.

\begin{remark}[Form and strength]
\label{rem:strength}
The bound uses $P(\kappa)$ (the Chebyshev theta sum),
not $\pi(\kappa)$.
The $P$-form is \emph{weaker} as a hypothesis than the $\pi$-form:
$P(\kappa)\sim\kappa$ while $\pi(\kappa)\sim\kappa/\log\kappa$,
so $P(\kappa)/\pi(\kappa)\sim\log\kappa\to\infty$;
the bound on $|M_k(\kappa)|$ with the $P$-normalisation has a larger right-hand side.
The $P$-form is nonetheless the numerically consistent choice:
$|M_k(\kappa)|/P(\kappa)$ stays bounded near $\pm 1/\gamma_k$
[numerical], while $|M_k(\kappa)|/\pi(\kappa)$ grows logarithmically [numerical].
Numerically, $|M_k(\kappa)|/P(\kappa)$ oscillates near $\pm 1/\gamma_k$
[numerical], consistent with a logarithmic saving of size $(\log\gamma_k)^{-1}$.
Assumption~\Astar\ asks only for a logarithmic saving in the
ordinate aspect, relative to the trivial $P(\kappa)$ bound.
It should not be read as a square-root-cancellation statement
in the cutoff variable $\kappa$; the present paper uses only the
displayed logarithmic-saving hypothesis.
The non-vanishing $\zeta(1+it)\ne 0$~\cite{Hadamard,dlVP} establishes
that there are no zeros on the line $\Re(s)=1$, but does not by itself
imply the logarithmic saving in $M_k(\kappa)$ required by
Assumption~\Astar.
The possible relation to Vinogradov--Korobov zero-free-region methods
\cite{MossinghoffTrudgianYang24,Bellotti24a}
remains open; see Open Problem~\ref{op:cancel}.
\end{remark}

\begin{remark}[Structural necessity]
\label{rem:necessity}
Five independent approaches to establishing
$|O'(\tfrac12)|=o(\Wgamma\cdot P(\kappa))$ without
Assumption~\Astar\ were investigated:
(i) Gaussian decay of the Guinand--Weil derivative as $\varepsilon\to 0$;
(ii) direct analysis of the cancellation mechanism via
     Definition~\ref{def:threeterm};
(iii) the functional equation $\zeta(s)=\chi(s)\zeta(1-s)$, which
     induces $\gamma_k\leftrightarrow-\gamma_k$ (reality of $O$) but does
     not induce the symmetry $O(\sigma)=O(1-\sigma)$;
(iv) three direct unconditional routes via Weyl equidistribution and
     trigonometric cancellation;
(v) the $\kappa$-dependence $O'(\tfrac12;101,0.05,100)=-137.6$, showing
     that smallness at the reference parameters is a numerical coincidence
     for this particular $\kappa$, not a universal structural property.
All five routes yielded only negative results or conditional statements.
These failures do not constitute a necessity proof, but they show that
the present paper has no unconditional route to stationarity;
Assumption~\Astar\ isolates the missing prime-exponential-sum input.
\end{remark}

% ============================================================
\section{The Integrated-Bias Window}
\label{sec:curv}
% ============================================================

\begin{proposition}[Integrated-bias window; proved under
                    Assumptions~\ref{ass:subleading} and~\ref{ass:proxy}]
\label{prop:intbias}
For every prime cutoff $\kappa\ge 2$ and every prime $p\le\kappa$,
$Z_{p,\infty}^+(\varepsilon)<0$ for all sufficiently small $\varepsilon>0$.
Consequently there exists
$\varepsilon_{\mathrm{int}}(\kappa)>0$ such that
\[
B_{\mathrm{int}}^\infty(\kappa,\varepsilon)
  := \sum_{p\le\kappa}(\log p)^2\,Z_{p,\infty}^+(\varepsilon) < 0
\]
for all $0<\varepsilon<\varepsilon_{\mathrm{int}}(\kappa)$.
\end{proposition}

\begin{proof}
Theorem~\ref{thm:rpinf} gives $r_p^\infty=\tfrac12<1$, and
Corollary~\ref{cor:sign} gives $Z_{p,\infty}^+(\varepsilon)<0$ for all
$\varepsilon<\varepsilon_0(p)$.
Since there are finitely many primes $p\le\kappa$, setting
$\varepsilon_{\mathrm{int}}(\kappa):=\min_{p\le\kappa}\varepsilon_0(p)>0$ gives
$Z_{p,\infty}^+(\varepsilon)<0$ for all $p\le\kappa$ and all
$\varepsilon\in(0,\varepsilon_{\mathrm{int}}(\kappa))$.
Each summand $(\log p)^2\,Z_{p,\infty}^+(\varepsilon)<0$, so the
sum $B_{\mathrm{int}}^\infty(\kappa,\varepsilon)<0$.
\end{proof}

\begin{remark}
The curvature identity $O''(\tfrac12)=-2B$ involves the
$\gamma_k^2$-weighted sum
\[
B \;:=\; \sum_p(\log p)^2\,\mathrm{Re}\bigl(Z_p^{(2)}\bigr),
\qquad
Z_p^{(2)} \;:=\; \sum_k\gamma_k^2\,e^{-\varepsilon^2\gamma_k^2}\,p^{i\gamma_k},
\]
which carries an additional spectral weight absent from
$B_{\mathrm{int}}^\infty$.
Whether $B_{\mathrm{int}}^\infty<0$ implies $B<0$ is the
\emph{Weighted Bias Bridge} (Paper~5, \S\,6;
Open Problem~\ref{op:wbridge}).
At reference parameters $(\kappa,\varepsilon,N)=(53,0.05,100)$,
$O''(\tfrac12)=+38\,685>0$ is verified numerically.
\end{remark}

\begin{remark}[Numerical values]
At the reference parameters $\varepsilon_{\mathrm{int}}(53)>0.200$; numerical
computation gives $\varepsilon_{\mathrm{int}}(2003)\approx 0.086$ and
$\varepsilon_{\mathrm{int}}(5003)\approx 0.069$ [numerical].
\end{remark}

\begin{remark}[Non-monotonicity]
The function $\kappa\mapsto\varepsilon_{\mathrm{int}}(\kappa)$ is non-monotone:
sign failures of $B_{\mathrm{int}}^\infty$ occur at isolated values
($\kappa=2003$, $5003$, $20011$) while $\varepsilon_{\mathrm{int}}(10007)>0.200$
[numerical].
This non-monotonicity is driven by a last-prime oscillation effect:
at large $\varepsilon$, where only $N_{\mathrm{eff}}=5$--$15$ zeros
contribute significantly, the sign of $B_{\mathrm{int}}^\infty(\kappa,\varepsilon)$
depends on the Gaussian phase of the final primes before the cutoff.
The existence of $\varepsilon_{\mathrm{int}}(\kappa)>0$ is conditional on
Assumption~\ref{ass:subleading},
as Proposition~\ref{prop:intbias} establishes.
\end{remark}

% ============================================================
\section{Conditional Selection Criterion}
\label{sec:sel}
% ============================================================

\begin{proposition}[Derivative bound; conditional under~\Astar]
\label{prop:deriv}
Assume Assumption~\Astar\ with constants $C_0>0$ and $A\ge 1$.
For all $\kappa\ge 2$, $\varepsilon>0$, and $N\ge 1$,
\[
\left|O'\!\left(\tfrac12;\kappa,\varepsilon,N\right)\right|
\;\le\; C\cdot\Wgamma\cdot P(\kappa),
\qquad C:=\frac{C_0}{(\log\gamma_1)^A},
\]
where $\log\gamma_1\approx 2.6486$ is independent of
$\kappa,\varepsilon,N$.
If additionally $C_0/(\log\gamma_1)^A<1$, this bound is
genuinely smaller than the trivial bound $\Wgamma\cdot P(\kappa)$.
\end{proposition}

\begin{proof}
By Assumption~\Astar, $|M_k(\kappa)|\le C_0\,P(\kappa)/(\log\gamma_k)^A$
for every $k\ge 1$.
Since $O'(\tfrac12)=-\sum_{k=1}^N e^{-\varepsilon^2\gamma_k^2}\gamma_k\,M_k(\kappa)$,
we obtain
\[
|O'(\tfrac12)|
  \;\le\; C_0\,P(\kappa)\sum_{k=1}^N
         \frac{e^{-\varepsilon^2\gamma_k^2}\,\gamma_k}{(\log\gamma_k)^A}
  \;\le\; \frac{C_0\,P(\kappa)}{(\log\gamma_1)^A}
         \sum_{k=1}^N e^{-\varepsilon^2\gamma_k^2}\,\gamma_k
  \;=\; C\cdot\Wgamma\cdot P(\kappa),
\]
where the second inequality uses $\log\gamma_k\ge\log\gamma_1>0$
for all $k\ge 1$.
\end{proof}

\begin{remark}
For fixed $\varepsilon>0$, the infinite analogue
$\sum_{k\ge 1}e^{-\varepsilon^2\gamma_k^2}\gamma_k$
converges by Gaussian decay together with the Riemann--von Mangoldt
zero-counting estimate $N(T)\sim T\log T/(2\pi)$
\cite[Thm.\,5.8]{IwaniecKowalski04}; the bound of
Proposition~\ref{prop:deriv} therefore also holds in the limit $N\to\infty$.
\end{remark}

\begin{corollary}[Integrated-bias input; conditional under
  Assumptions~\ref{ass:subleading} and~\ref{ass:proxy}]
\label{cor:intbias_combined}
Under Assumptions~\ref{ass:subleading} and~\ref{ass:proxy},
$B_{\mathrm{int}}^\infty(\kappa,\varepsilon)<0$ for
$\varepsilon\in(0,\varepsilon_{\mathrm{int}}(\kappa))$
(Proposition~\ref{prop:intbias}).
\end{corollary}

\begin{corollary}[Second-order local selection criterion]
\label{cor:selection}
Fix $(\kappa,\varepsilon,N)$. If $B(\kappa,\varepsilon,N)<0$ and
$O'(\tfrac12;\kappa,\varepsilon,N)=0$, then $\sigma=\tfrac12$ is a
strict local minimum of
$\sigma\mapsto O(\sigma;\kappa,\varepsilon,N)$, and
$\mathrm{Tr}(\Tt(\sigma))$ has a strict local maximum
at $\sigma=\tfrac12$.
\end{corollary}

\begin{proof}
Under the stated hypotheses, $O'(\tfrac12)=0$ and the curvature
identity $O''(\tfrac12)=-2B$ of~\cite{Paper6} gives
$O''(\tfrac12)>0$.
The second-order condition then gives a strict local minimum,
hence a strict local maximum of
$\mathrm{Tr}(\Tt(\sigma))=\DSEL-O(\sigma)$.
\end{proof}

The conditional inputs of Proposition~\ref{prop:deriv}
(Assumption~\Astar) and Corollary~\ref{cor:intbias_combined}
(Assumptions~\ref{ass:subleading} and~\ref{ass:proxy}) do not
directly imply $B<0$ or $O'(\tfrac12)=0$: supplying $B<0$
requires the Weighted Bias Bridge (Open Problem~\ref{op:wbridge}),
and exact stationarity requires a separate analytical argument.
The conditional hierarchy of \S\,\ref{sec:sel} therefore stands
as follows:
\[
\text{Assumptions~\ref{ass:subleading} and~\ref{ass:proxy}}
  \;\Longrightarrow\; B_{\mathrm{int}}^\infty<0
  \qquad\text{(Cor.~\ref{cor:intbias_combined})};
\]
\[
\Astar \;\Longrightarrow\; |O'(\tfrac12)|\le C\cdot\Wgamma\cdot P(\kappa)
  \qquad\text{(Prop.~\ref{prop:deriv})};
\]
\[
B<0 \text{ and } O'(\tfrac12)=0
  \;\Longrightarrow\; \text{strict local minimum at }\sigma=\tfrac12
  \qquad\text{(Cor.~\ref{cor:selection})}.
\]
At the reference parameters $(\kappa,\varepsilon,N)=(53,0.05,100)$,
$B=-19\,342.5<0$ is verified numerically; for general parameters,
$B<0$ is the content of Open Problem~\ref{op:wbridge}.

% ============================================================
\section{Structural Remarks}
\label{sec:struct}
% ============================================================

\subsection{Antisymmetry of the spectral response}

\begin{proposition}[Antisymmetry; proved]
\label{prop:antisym}
Define $\mathcal{A}(\sigma):=O(\sigma)-O(1-\sigma)$.
For all $\sigma\in\mathbb{R}$ and $\delta\in\mathbb{R}$,
\[
\mathcal{A}(\sigma)
  \;=\; -\sum_{k,p}e^{-\varepsilon^2\gamma_k^2}
    \sin(\gamma_k\log p)\,\sin\bigl((2\sigma-1)\gamma_k\log p\bigr),
\]
and in particular $\mathcal{A}(\tfrac12+\delta)=-\mathcal{A}(\tfrac12-\delta)$.
\end{proposition}

\begin{proof}
Apply the difference formula $\cos\alpha-\cos\beta
=-2\sin\tfrac{\alpha+\beta}{2}\sin\tfrac{\alpha-\beta}{2}$
to each term of $O(\sigma)-O(1-\sigma)$, with
$\alpha=2\sigma\gamma_k\log p$ and $\beta=2(1-\sigma)\gamma_k\log p$.
One obtains $\alpha-\beta=2(2\sigma-1)\gamma_k\log p$ and
$(\alpha+\beta)/2=\gamma_k\log p$, yielding the stated formula.
The reflection $\sigma\mapsto 1-\sigma$ sends $2\sigma-1\mapsto-(2\sigma-1)$,
so $\mathcal{A}(1-\sigma)=-\mathcal{A}(\sigma)$; setting $\sigma=\tfrac12+\delta$
gives $\mathcal{A}(\tfrac12+\delta)=-\mathcal{A}(\tfrac12-\delta)$.
\end{proof}

\begin{remark}
The symmetry $O(\sigma)=O(1-\sigma)$ does not hold:
one verifies algebraically that $O'(\tfrac12)\ne 0$ in general
(numerically, $O'(\tfrac12;53,0.05,100)=+2.4751\ne 0$).
The functional equation $\zeta(s)=\chi(s)\zeta(1-s)$ induces a
reflection symmetry $\gamma_k\leftrightarrow-\gamma_k$ (reality of
$O$), but not $\sigma\leftrightarrow 1-\sigma$.
\end{remark}

\subsection{Remark on unconditional routes to stationarity}

Investigation of five independent approaches to establishing
$|O'(\tfrac12)|=o(\Wgamma\cdot P(\kappa))$ without
Assumption~\Astar\ yielded only negative results or conditional statements
(see Remark~\ref{rem:necessity} for details).
In particular, the $\kappa$-dependence of $O'(\tfrac12)$ --- which equals
$-137.6$ at $\kappa=101$ for the same $(\varepsilon,N)$ --- shows that the
smallness at $\kappa=53$ is a numerical coincidence, not a universal
structural identity.
The deep cancellation witnessed numerically has the character of
a trigonometric sum estimate, and the Cancellation Hypothesis~\Astar\
captures precisely the arithmetic content required.

% ============================================================
\section{Open Problems}
\label{sec:open}
% ============================================================

\begin{openproblem}[Weighted Bias Bridge]
\label{op:wbridge}
\emph{Given} (Proposition~\ref{prop:intbias}, under
Assumptions~\ref{ass:subleading} and~\ref{ass:proxy}):
\[
B_{\mathrm{int}}^\infty(\kappa,\varepsilon)
 := \sum_p(\log p)^2\, Z_{p,\infty}^+(\varepsilon)
 \;<\; 0
\qquad
\text{for } 0<\varepsilon<\varepsilon_{\mathrm{int}}(\kappa).
\]

\emph{Target}: the direct curvature sum
\[
B(\kappa,\varepsilon,N)
 := \sum_p(\log p)^2\,\mathrm{Re}\bigl(Z_p^{(2)}\bigr),
\]
which gives $O''(\tfrac12)=-2B$ via~\cite{Paper6}.

The bridge requires two successive transfers:

\smallskip
\noindent\emph{Transfer (i): unweighted $\to\,\gamma_k^2$-weighted asymptotic.}
Pass from $B_{\mathrm{int}}^\infty$ to
\[
B^{(2),\infty}(\kappa,\varepsilon)
 := \sum_p(\log p)^2\,\mathrm{Re}\!\left(
    \sum_{k\ge 1}\gamma_k^2\,e^{-\varepsilon^2\gamma_k^2}
    \,p^{i\gamma_k}
    \right).
\]
This requires controlling the $\gamma_k^2$-weighting in the
infinite ordinate sum.

\smallskip
\noindent\emph{Transfer (ii): asymptotic $\to$ finite-$N$.}
Pass from $B^{(2),\infty}(\kappa,\varepsilon)$ to
$B(\kappa,\varepsilon,N)$.
This requires a finite-$N$ truncation argument.

\smallskip
Together, Transfers~(i) and~(ii) would yield $B<0$, hence
$O''(\tfrac12)>0$ at the reference parameters.
This is the \emph{Weighted Bias Bridge} of Paper~5, \S\,6.
\end{openproblem}

\begin{openproblem}[Cancellation Hypothesis]
\label{op:cancel}
Establish Assumption~\Astar\ analytically.
In particular: does there exist a treatment via Vinogradov--Korobov
exponential sum techniques
(see \cite{MossinghoffTrudgianYang24,Bellotti24a} for the
state of the art on zero-free regions) which yields the logarithmic
saving in $|M_k(\kappa)|/P(\kappa)$ asserted by Assumption~\Astar?
\end{openproblem}

\begin{openproblem}[GW bridge and off-critical control]
\label{op:gwbridge}
The full Guinand--Weil zero sum satisfies
\[
\widetilde Z_p^{\mathrm{GW,full}}(\varepsilon)
= Z_p^\infty(\varepsilon) + E_p^{\mathrm{off}}(\varepsilon),
\]
where $E_p^{\mathrm{off}}$ collects contributions from zeros
off the critical line (Paper~5, \S\,6).
Establish bounds or sign control on $E_p^{\mathrm{off}}$,
or identify conditions under which
$E_p^{\mathrm{off}}(\varepsilon)=o(|\mathrm{Main}_p(\varepsilon)|)$
holds without assuming RH.
(This is the same scale as Assumption~\ref{ass:proxy};
under Theorem~\ref{thm:rpinf} the two scales
$|\mathrm{Main}_p|$ and $|Z_p^\infty|$ are comparable.)
The present paper proves $r_p^\infty=\tfrac12$ for $Z_{p,\infty}^+$
(under Assumptions~\ref{ass:subleading} and~\ref{ass:proxy});
the unconditional transfer to $\widetilde Z_p^{\mathrm{GW,full}}$
requires controlling $E_p^{\mathrm{off}}$.
\end{openproblem}

\begin{openproblem}[Scaling of $\varepsilon_{\mathrm{int}}(\kappa)$]
\label{op:epsc}
Determine the asymptotic behaviour of $\varepsilon_{\mathrm{int}}(\kappa)$ as $\kappa\to\infty$.
In particular: is $\inf_\kappa\varepsilon_{\mathrm{int}}(\kappa)>0$?
Numerical evidence gives $\varepsilon_{\mathrm{int}}(53)>0.200$, $\varepsilon_{\mathrm{int}}(2003)\approx 0.086$,
$\varepsilon_{\mathrm{int}}(5003)\approx 0.069$, $\varepsilon_{\mathrm{int}}(10007)>0.200$ [numerical],
suggesting oscillatory behaviour without monotone decrease, but no
analytical lower bound is known.
\end{openproblem}

\begin{openproblem}[Near-minimum behaviour at the critical line]
\label{op:global}
For which parameter pairs $(\kappa,\varepsilon,N)$ does the continuous
minimum of $O(\sigma;\kappa,\varepsilon,N)$ lie near $\sigma=\tfrac12$,
and under what additional stationarity condition can it occur exactly
at $\sigma=\tfrac12$?
At the reference parameters $(53,0.05,100)$ one has
$O'(\tfrac12)=+2.4751\ne 0$, so $\sigma=\tfrac12$ is not an exact
continuous local or global minimum.
The nearby stationary point lies at
$\sigma_*\approx 0.4999$ [numerical, resolved local scale].
Exact selection of $\sigma=\tfrac12$ requires the separate stationarity
condition isolated in Corollary~\ref{cor:selection}.
For $\kappa=101$, $\varepsilon=0.05$ (fixed $N=100$), the corresponding
near-minimum behaviour fails [numerical].
\end{openproblem}

\begin{openproblem}[Weil-transfer operator and positivity]
\label{op:weil}
Paper~2~\cite{Paper2} introduces $\sigma$-dependent admissible Gaussian
test functions $g^*_{\sigma,\varepsilon}\in\mathcal{S}_{\mathrm{ad}}$,
but does not provide a per-prime family.
For the present problem, define the prime-local antisymmetric
Gaussian component
\[
\psi_{p,\varepsilon}(x)
  := \phi_\varepsilon(x-\log p) - \phi_\varepsilon(x+\log p),
\]
where $\phi_\varepsilon$ is the Gaussian
$\phi_\varepsilon(x)=(2\pi\varepsilon^2)^{-1/2}\exp(-x^2/(2\varepsilon^2))$.
Define the \emph{Weil-transfer matrix}
$\mathsf{W}^{\mathrm{Weil}}_{\kappa,\varepsilon}$ by
\[
(\mathsf{W}^{\mathrm{Weil}}_{\kappa,\varepsilon})_{pq}
  \;:=\; W\!\left(\psi_{p,\varepsilon}\ast\check\psi_{q,\varepsilon}\right),
\]
where $W(\cdot)$ denotes the Weil quadratic functional of
\cite{Lagarias,Bombieri}.
The matrix $\mathsf{W}^{\mathrm{Weil}}_{\kappa,\varepsilon}$ is a finite-dimensional
representation of the Weil form on the admissible subspace
$H_{\kappa,\varepsilon}^{\mathrm{test}}
:=\mathrm{span}\{\psi_{p,\varepsilon}:p\le\kappa\}$.

\emph{(i)} Establish whether the bridge identity holds,
where $c^\eta:=(\sqrt{V_p(\tfrac12)})_{p\le\kappa}$ is the vector
of spectral weights introduced in Paper~3~\cite{Paper3}:
\[
\bigl\langle\mathsf{W}^{\mathrm{Weil}}_{\kappa,\varepsilon}\,c^\eta,\,c^\eta\bigr\rangle
  \;=\; \alpha_{\kappa,\varepsilon}\cdot O''\!\left(\tfrac12\right)
    + E_{\mathrm{off}} + E_{\Gamma/\mathrm{pole}}
\]
holds for some normalisation constant $\alpha_{\kappa,\varepsilon}$,
with $E_{\mathrm{off}}$ an off-diagonal remainder and
$E_{\Gamma/\mathrm{pole}}$ a $\Gamma$-factor correction.
The four open verification steps are:
admissibility of the family $\{\psi_{p,\varepsilon}\}$ in the Weil
sense and its compatibility with the Paper-2 normalisation and
Fourier convention;
derivation of $\alpha_{\kappa,\varepsilon}$ from Fourier conventions;
formal identification of the diagonal
$(\mathsf{W}^{\mathrm{Weil}}_{\kappa,\varepsilon})_{pp}\sim V_p(\tfrac12)$;
and control of $E_{\mathrm{off}}$.

\emph{(ii)} If~(i) holds:
can the conditional input $B_{\mathrm{int}}^\infty<0$ from
Proposition~\ref{prop:intbias}, together with an independent proof
or finite-grid verification of $O''(\tfrac12)>0$, be transferred to
a Weil-functional shadow
$\langle\mathsf{W}^{\mathrm{Weil}}_{\kappa,\varepsilon}\,c^\eta,c^\eta\rangle\ge 0$
on a growing finite-cutoff subspace, and does
$\mathsf{W}^{\mathrm{Weil}}_{\kappa,\varepsilon}\succeq 0$ hold in the limit?
Positivity of the Weil form on this subspace would constitute a
finite-cutoff analogue of the Weil positivity criterion of
\cite{Lagarias}, which is known to be equivalent to the
Riemann Hypothesis~\cite{Lagarias,Bombieri}.
\end{openproblem}

% ============================================================
\section*{Non-Circularity}
\label{sec:nocirc}
% ============================================================

We document explicitly that the results of this paper do not
rely on the objects they are directed toward.

\textbf{No RH as input.}
At no point is the Riemann Hypothesis assumed.
The ordinates $\gamma_k$ are used as numerically specified inputs;
no property of their distribution requiring the Riemann Hypothesis
is used in any proof.
The identity $\widetilde Z_p^{\mathrm{GW,full}}=Z_p^\infty$ (under RH)
is carefully separated from the unconditional computations, which use
only the ordinate proxy $Z_{p,N}^+$.

\textbf{No GUE, Montgomery, or Hilbert--P\'olya hypothesis.}
The distribution of the zeros is not modelled by random matrix theory.
The operator $\Tt$ is not a Hilbert--P\'olya operator (Paper~4,
Hilbert--P\'olya diagnostic: $r_2\to 0.16$); its matrix entries are
explicit deterministic arithmetic expressions.
No pair correlation conjecture or related statistical hypothesis is used.

\textbf{Conditional inputs.}
Theorem~\ref{thm:rpinf}, Corollary~\ref{cor:sign}, and
Proposition~\ref{prop:intbias} are conditional on
Assumptions~\ref{ass:subleading} and~\ref{ass:proxy}.
Proposition~\ref{prop:deriv} is conditional on
Assumption~\Astar\ (Cancellation Hypothesis).
Definition~\ref{def:threeterm}, Observation~\ref{obs:threeterm},
and Proposition~\ref{prop:antisym} are unconditional.
Neither Assumption~\ref{ass:subleading} nor Assumption~\Astar
is known to follow from RH, GUE, or a Hilbert--P\'{o}lya postulate.
Assumption~\ref{ass:proxy} is implied by RH and is formally weaker
in content: it requires only asymptotic smallness
$o(|\mathrm{Main}_p|)$ of the off-critical contribution, not its
vanishing. No converse implication is claimed.
The non-vanishing $\zeta(1+it)\ne 0$~\cite{Hadamard,dlVP}
is an unconditional analytical input, distinct from both assumptions.

% ============================================================
\section*{Acknowledgments}
% ============================================================

The author thanks the readers who provided feedback on earlier
versions of this paper and on the preceding papers in the series.

% ============================================================
\section*{Call for Feedback}
% ============================================================

This paper is part of an ongoing open research initiative toward
the Riemann Hypothesis.
All papers in the series are freely available on Zenodo under CC~BY~4.0,
and the accompanying verification code is hosted on GitHub.

{\sloppy
Zenodo (\url{https://zenodo.org/search?q=Ulrich+Tehrani}) and GitHub
(\url{https://github.com/utehrani/analysislab-nt}).
\par}

Topics of particular interest include:
analytical approaches to Assumption~\Astar\ via exponential sum methods
and Vinogradov--Korobov zero-free region techniques;
connections between the integrated-bias window $\varepsilon_{\mathrm{int}}(\kappa)$
and zero-free regions for $\zeta(s)$;
and any errors or gaps in the present arguments.

% ============================================================
\newpage
\begin{thebibliography}{99}

\bibitem{Paper6}
U.~Tehrani,
``Positive Curvature of the Spectral Trace at the Critical Line'',
Zenodo, \url{https://doi.org/10.5281/zenodo.19665790}, 2026.

\bibitem{Paper5}
U.~Tehrani,
``Spectral Trace Formula and Smoothed Zero Sums:
A Prime--Zero Duality Framework'',
Zenodo, \url{https://doi.org/10.5281/zenodo.19508547}, 2026.

\bibitem{Code}
U.~Tehrani,
Verification scripts for the prime--zero coupling series,
GitHub, \url{https://github.com/utehrani/analysislab-nt}.

\bibitem{Paper1}
U.~Tehrani,
``A Curvature Decomposition of the Explicit Formula
for the Riemann Zeta Function'',
Zenodo, \url{https://doi.org/10.5281/zenodo.19025598}, 2026.

\bibitem{Paper2}
U.~Tehrani,
``From Local Curvature to the Weil Functional:
An Explicit Construction'',
Zenodo, \url{https://doi.org/10.5281/zenodo.19106992}, 2026.

\bibitem{Lagarias}
J.C.~Lagarias,
``On a positivity property of the Riemann $\xi$-function'',
\emph{Acta Arithmetica} \textbf{89} (1999), 217--234.

\bibitem{Paper3}
U.~Tehrani,
``A Finite-Cutoff Hilbert-Space Model for Prime--Zero Energy Structure'',
Zenodo, \url{https://doi.org/10.5281/zenodo.19307989}, 2026.

\bibitem{Paper4}
U.~Tehrani,
``A Dual Operator for Prime--Zero Coupling
and a Conditional Proof of Energy Asymmetry'',
Zenodo, \url{https://doi.org/10.5281/zenodo.19364703}, 2026.

\bibitem{BerryKeating}
M.V.~Berry and J.P.~Keating,
``The Riemann Zeros and Eigenvalue Asymptotics'',
\emph{SIAM Review} \textbf{41}.2 (1999), 236--266.

\bibitem{deBranges}
L.~de~Branges,
``The Riemann Hypothesis for Hilbert Spaces of Entire Functions'',
\emph{Bull.\ Amer.\ Math.\ Soc.\ (N.S.)} \textbf{15}.1 (1986), 1--17.

\bibitem{ConnesConsani23}
A.~Connes and C.~Consani,
``Spectral triples and zeta-cycles,''
\textit{Enseign.\ Math.} \textbf{69} (2023), no.~1--2, 93--148.
arXiv:2106.01715.

\bibitem{CCM25}
A.~Connes, C.~Consani, and H.~Moscovici,
``Zeta spectral triples,''
arXiv:2511.22755 (2025).

\bibitem{Bombieri}
E.~Bombieri,
``Remarks on Weil's Quadratic Functional in the Theory of Prime
Numbers, I'',
\emph{Rend.\ Lincei Mat.\ Appl.} \textbf{11}.3 (2000), 183--233.

\bibitem{Burnol}
J.-F.~Burnol,
``On Fourier and Zeta(s)'',
\emph{Forum Math.} \textbf{16}.6 (2004), 789--840.

\bibitem{Suzuki25}
M.~Suzuki,
``On the Hilbert space derived from the Weil distribution,''
\textit{Canad.\ J.\ Math.}, First View (2025), 1--24.
DOI: 10.4153/S0008414X25101739.
arXiv:2301.00421.

\bibitem{Guinand}
A.P.~Guinand,
``A summation formula in the theory of prime numbers'',
\emph{Proc.\ London Math.\ Soc.}\ (2) \textbf{50} (1948), 107--119.

\bibitem{Weil}
A.~Weil,
``Sur les `formules explicites' de la th\'{e}orie des nombres premiers'',
\emph{Comm.\ S\'{e}m.\ Math.\ Univ.\ Lund, Suppl.}\ (1952), 252--265.

\bibitem{IwaniecKowalski04}
H.~Iwaniec and E.~Kowalski,
\emph{Analytic Number Theory},
American Mathematical Society Colloquium Publications,
vol.~53, AMS, Providence, RI, 2004.

\bibitem{MossinghoffTrudgianYang24}
M.~J.~Mossinghoff, T.~S.~Trudgian, and A.~Yang,
``Explicit zero-free regions for the Riemann zeta-function,''
\textit{Res.\ Number Theory} \textbf{10} (2024), art.~11.
DOI: 10.1007/s40993-023-00498-y.

\bibitem{Bellotti24a}
C.~Bellotti,
``Explicit bounds for the Riemann zeta function
and a new zero-free region,''
\textit{J.\ Math.\ Anal.\ Appl.} \textbf{536} (2024),
no.~2, art.~128249.
DOI: 10.1016/j.jmaa.2024.128249.
arXiv:2306.10680.

\bibitem{Hadamard}
J.~Hadamard,
``Sur la distribution des z\'{e}ros de la fonction $\zeta(s)$
et ses cons\'{e}quences arithm\'{e}tiques'',
\emph{Bull.\ Soc.\ Math.\ France} \textbf{24} (1896), 199--220.

\bibitem{dlVP}
C.-J.~de~la~Vall\'{e}e~Poussin,
``Recherches analytiques sur la th\'{e}orie des nombres premiers'',
\emph{Ann.\ Soc.\ Sci.\ Bruxelles} \textbf{20} (1896), 183--256.

\end{thebibliography}
% ============================================================

\enlargethispage{3\baselineskip}
\vspace*{\fill}
\begin{minipage}{\linewidth}
\rule{\linewidth}{0.4pt}\smallskip\\
\small\emph{Paper~7 v0.14 \textperiodcentered\ May~2026
\textperiodcentered\ Pauca sed matura}
\end{minipage}

\end{document}
